% !TeX root = ../../Book5.tex
% 纲要标题：### 17.2 Cartan 矩阵与 Dynkin 图
% 纲要位置：工作记录/计划纲要.md，第 4264 行。

\section{Cartan 矩阵与 Dynkin 图}
\label{b5:sec:1702}

简单根数目有限，但其内积最初看似仍有连续参数。结晶整数把角度和长度比限制为少数可能，Cartan 矩阵由此成为一份离散的记录。

% 纲要要点已在下文展开。

% 纲要细目（保留原定写作范围）：
%
% * 简单根夹角和长度
% * Cartan 矩阵
% * Dynkin 图的构造
%

\subsection{简单根夹角和长度}
\label{b5:subsec:1702:01}

两个不成比例的根 \(\alpha,\beta\) 给出的整数
\[
a=\frac{2(\beta,\alpha)}{(\alpha,\alpha)},\qquad
b=\frac{2(\alpha,\beta)}{(\beta,\beta)}
\]
满足 \(ab=4\cos^2\theta<4\)。对不同简单根，\(a,b\leq0\)，所以可能的有序数对只有
\[
(0,0),\quad(-1,-1),\quad(-1,-2),\quad(-2,-1),
\quad(-1,-3),\quad(-3,-1).
\]
它们对应夹角 \(90^\circ,120^\circ,135^\circ,150^\circ\)；
非等长时，平方长度之比为 \(2\) 或 \(3\)。
这说明结晶条件不仅限制角度，也限制不同长度的相对比例。

\subsection{Cartan 矩阵}
\label{b5:subsec:1702:02}

\begin{definition}\label{b5:def:root-cartan-matrix}
设 \(\Delta=\{\alpha_1,\ldots,\alpha_r\}\) 为有限约化结晶根系的简单根基。矩阵
\[
A=(a_{ij}),\qquad
a_{ij}=\langle\alpha_j,\alpha_i^\vee\rangle
=\frac{2(\alpha_i,\alpha_j)}{(\alpha_i,\alpha_i)}
\]
称为根系的\term{Cartan 矩阵}[Cartan matrix]。本书的行对应余根，列对应根。
\end{definition}

这个矩阵与模表示论中由投射覆盖重数定义的 Cartan 矩阵是不同对象。这里
\(a_{ii}=2\)，非对角元为非正整数，且
\(a_{ij}=0\) 当且仅当 \(a_{ji}=0\)。
令 \(d_i=(\alpha_i,\alpha_i)/2\)，则
\(\operatorname{diag}(d_i)A\) 是简单根的 Gram 矩阵，故对称正定。特别地 \(A\) 可逆。

\begin{proposition}\label{b5:prop:cartan-determines-roots}
设两个有限约化结晶根系具有相同的带编号 Cartan 矩阵。把一组简单根送到另一组简单根的线性同构，将全部根双射到全部根；在每个不可约分支上，它将内积乘以一个正数。
\end{proposition}
\begin{proof}
在简单根基下，
\(s_i(\alpha_j)=\alpha_j-a_{ij}\alpha_i\)，所以这个线性同构交织全部简单反射。
每个根都是某个简单根经这些反射所得，故它将根集合双射到根集合。
在相邻顶点处，\(a_{ij}/a_{ji}\) 决定平方长度之比；沿连通图传播，所有长度只差一个共同正比例。再由 Cartan 系数恢复全部内积，便得到最后的断言。
\end{proof}

\subsection{Dynkin 图的构造}
\label{b5:subsec:1702:03}

为每个简单根画一个顶点；顶点 \(i,j\) 之间画 \(a_{ij}a_{ji}\) 条边，若是双边或三边，则箭头指向较短的根。所得有向重图称为 Dynkin 图。它同时记录角度及长短，因而恢复 Cartan 矩阵。例如
\[
A_2:\ \begin{pmatrix}2&-1\\-1&2\end{pmatrix},\qquad
B_2:\ \begin{pmatrix}2&-1\\-2&2\end{pmatrix},\qquad
G_2:\ \begin{pmatrix}2&-3\\-1&2\end{pmatrix}.
\]
在这里的 \(B_2\) 编号中第一根长、第二根短；\(G_2\) 的第一根短。
若只保留边数而忘掉箭头，就不能区分一般秩的 \(B_r\) 与 \(C_r\)。

\ProseParagraphBreak

分类时方便把每根归一到长度 \(\sqrt2\)。其 Gram 矩阵 \(G\) 满足
\[
G_{ii}=2,\qquad G_{ij}=-\sqrt{a_{ij}a_{ji}}\quad(i\ne j).
\]
这只是对 Gram 矩阵作正对角合同变换，仍然正定；它不把原来的长短根真的改成同一个根系。下一节通过这个正定矩阵限制图的形状，再用具体坐标验证所有剩下的图确实出现。
